\chapter{Discussion and Conclusion}
\label{ch:disc}

\chapterhandoff{Chapters~\ref{ch:repr}--\ref{ch:e2e} established a ceiling, a
failure to predict, and a failure to profit. This chapter explains the second of
those in terms of a mechanism, states the regime in which all three hold, and
answers the research question at exactly the level of evidence available.}

\section{Why the prediction failed: a supply constraint}
\label{sec:mechanism}

The natural explanations for a failed classifier are the hypothesis class, the
feature set, or the label definition. On this data none of them is the binding
constraint. The binding constraint is that \textbf{routing labels are produced
only by solve events, and solve events are scarce}.

Consider the attrition (Figure~\ref{fig:f14}). On \cell{cls$\cdot$B0}, the most
favourable cell in the study, 224 scored tasks yield 97 tasks that some mode
solves --- and therefore 97 which-mode labels. Of the six possible classes, only
three retain enough members to survive a five-fold split. On
\cell{red$\cdot$B2}, 203 tasks yield 15 labels and no class survives at all.
Pooled over the six VisualWebArena cells there are 260 which-mode labels in
total.

\begin{figure}[t]
  \centering
  \includegraphics[width=\textwidth]{fig_f14_label_supply_attrition.pdf}
  \caption[Label-supply attrition per cell]{%
  \textbf{From scored tasks to trainable classes.} Each cell's funnel: scored
  tasks $\to$ tasks any mode solves $\to$ which-mode labels $\to$ classes
  retaining enough members to survive a five-fold split. The horizontal line is
  the two-class minimum below which nothing can be discriminated. Four of six
  cells finish below it. The bottleneck is the second arrow, and its width is
  the base success rate --- not a modelling choice.}
  \label{fig:f14}
\end{figure}

This is a \emph{rate} problem, not a \emph{slicing} problem. No re-partition of
the supervision creates a solve event that the benchmark did not produce. Three
independent routes out of the constraint are available in principle; we
attempted all three, and each is blocked for a structurally different reason
(Figure~\ref{fig:f15}).

\subsection{Route 1 --- make the label continuous}

If success were graded rather than binary, every task would carry signal, and
scarcity would disappear. VisualWebArena's evaluator returns a binary verdict.
Its component evaluators (URL match, program-state HTML checks, string match,
image query) each emit a pass or fail that is combined conjunctively; there is
no partial-credit channel to expose. The route is closed by the benchmark's
measurement instrument, not by our analysis.

\subsection{Route 2 --- pool across cells}

Pooling tasks across backbones multiplies the label count. It also destroys
identifiability, because the same task acquires different labels under different
backbones. On classifieds, 54 tasks are covered by two or more cells and
\textbf{31 of them (57.4\%) carry conflicting labels}; on reddit the conflict
rate is 56.0\% of 25 shared tasks. A pooled classifier is being asked to predict
a quantity that is not a function of its inputs.

We report the associated agreement ceilings with an explicit caveat that they
are \emph{not} what they may look like. Grouping by task and predicting the
modal label attains 79.2\% (classifieds) and 83.7\% (reddit); grouping by exact
feature vector attains more. Both are \emph{resubstitution} estimates --- each
scores the same rows it took its modal label from, so any singleton group is
correct by construction, and singletons are 29.8\% and 37.0\% of rows
respectively. They are in-sample modal agreement, not Bayes ceilings, and we do
not use them as bounds.

\subsection{Route 3 --- relabel to a coarser target}

Re-slicing the same features from ``which of six modes'' to ``which cost tier''
raises the attainable agreement (89.9\% on classifieds, 96.7\% on reddit)
without inventing a single new solve event. It is the only relabelling that buys
anything, and what it buys is a coarser decision, not more supervision. The
route is open but leads somewhere smaller than where we were going.

\begin{figure}[t]
  \centering
  \includegraphics[width=\textwidth]{fig_f15_three_routes.pdf}
  \caption[Three routes out of the supply constraint, and why each is blocked]{%
  \textbf{Each escape from label scarcity fails differently.} Continuous labels
  are blocked by the evaluator, which emits a conjunctive binary verdict.
  Pooling is blocked by identifiability: 56--57\% of tasks shared by two or more
  cells carry conflicting labels across backbones. (The bars group by task
  identity; because three features are observation-derived, the same task can
  also carry different feature vectors across cells, and that stricter
  feature-space version is a separate diagnostic in the source artefact.) Coarser relabelling is open but converts
  the problem into a smaller one rather than supplying more evidence. Three
  distinct failure modes is itself the argument that the constraint is
  structural.}
  \label{fig:f15}
\end{figure}

\subsection{A defect we found in our own supervision}

Auditing the label revealed a flaw worth recording, because it is a plausible
mistake for anyone building this pipeline. Our canonical which-mode label is
``the first successful mode in a list ordered by prior cost''. The list is
documented as ascending in cost. It is not, against measured cost: on
12.5--54.6\% of labels, depending on cell, the ordering returns a mode that is
strictly \emph{more} expensive than another successful one. The label is
therefore not ``cheapest successful mode''; it is whatever the list reached
first.

We had earlier recorded this as a tie-breaking problem affecting roughly a
quarter of labels. That was the wrong frame and the wrong quantity. Exact ties
occur in \textbf{zero} labels in every cell --- measured cost is a continuous
float, and two modes costing exactly the same does not happen. The real defect
is not arbitrary tie-breaking but a false ordering assumption, which is both
more common and worse. Table~\ref{tab:supply} and its sensitivity analysis in
Section~\ref{sec:robustness} are computed under both labels, and the
supply conclusion is unchanged.

\section{Oracle headroom is not learnable headroom}
\label{sec:headroom-gap}

Chapter~\ref{ch:repr} measured a ceiling of $+3.45$ to $+16.35\pp$;
Chapter~\ref{ch:e2e} realised none of it. The gap decomposes into four terms,
each of which we can now attach a measurement to:
\begin{equation}
  \underbrace{\text{realised value}}_{\approx 0}
  \;\approx\;
  \underbrace{\text{selection opportunity}}_{+3.45\ \text{to}\ +16.35\pp}
  \times
  \underbrace{\text{decision quality}}_{\text{AUROC } 0.53\text{--}0.76}
  \;-\;
  \underbrace{\text{routing overhead}}_{\sim10^{-6}\,\$}
  \;-\;
  \underbrace{\text{error cost}}_{\text{dominant}} .
  \label{eq:decomposition}
\end{equation}

Equation~\ref{eq:decomposition} is a conceptual decomposition; we do not
estimate every term jointly. But the ordering of magnitudes is measured, and it
locates the loss unambiguously. The overhead term is negligible
(Section~\ref{sec:overhead}). The opportunity term is real. The loss is split
between decision quality --- weak, and mostly borrowed from a benchmark
annotation --- and error cost, which in a regime where 48--93\% of tasks are
hopeless is dominated by false escalations on tasks that were never going to
succeed.

There is also a fifth term, absent from the usual formulation, that this study
was forced to add: \textbf{the opportunity itself is largely repetition}. Adding
one \emph{rerun} arm buys 2.0--7.6\pp{} on its own
(Section~\ref{sec:noise-floor}); at the one-arm margin a distinct representation
is worth no more than a rerun (Figure~\ref{fig:f10b}); and at the matched arm
count of six, repetition of the strongest single mode closes to within 3.53\pp{}
of the six-mode oracle (Section~\ref{sec:six-vs-six}). An accounting that treats
the full ceiling as representation-attributable is therefore claiming something
our own data contradicts. The term that survives is not selection opportunity in
success but selection opportunity \emph{per episode spent} --- which is the axis
on which routing was always supposed to pay, and which makes the negative result
of Chapter~\ref{ch:e2e} sharper rather than softer.

\section{What the phantom-space result does and does not mean}
\label{sec:phantom-meaning}

The four image-free modes are a clean experimental factorisation: text payload
and prompt framing vary independently with the image held absent, and the
$1.008\times$ text ratio between \mDOM{} and \mSoM{} means the manipulation is
surgical. As an instrument, it works.

As a finding, it is weaker than we initially read it. The three phantom modes
contribute in a consistent direction --- 14 of 18 arm--cell drop-one values
positive --- but at pooled two-axis magnitudes of 1.35 and 2.09\pp, below even
the most permissive measured repetition floor. We therefore report a
directionally consistent complementarity structure and decline the structural
claim (Section~\ref{sec:dropone}).

It is fair to ask why a declined claim still gets a section, a table and a
figure. Two reasons, and neither is attachment to the result. The 36 drop-one
values are the evidence \emph{for} the decline --- deleting them would leave the
reader with an assertion where there is now a measurement, and a later reader
with no way to check whether the effect was small or absent. And the
factorisation that produced them is load-bearing for
Chapter~\ref{ch:system} regardless of what it found: it is what licenses
attributing anything at all to the image rather than to the text or the prompt.
What we do not do is let the structure carry weight elsewhere, and
Chapter~\ref{ch:e2e}'s result depends on none of it.

We also decline two adjacent readings that the data does not support. The
image-free modes are not a cheaper Set-of-Marks: \mSoM{} is the strongest single
mode in four of six VisualWebArena cells. And the modes are not evidence of
different internal reasoning mechanisms; they are evidence of different
input configurations producing different outcome sets. A mechanistic reading
would require separate evidence, which we did not collect.

\section{Sustainability: what was measured, and what follows}
\label{sec:sustainability-disc}

We measured GPU-device operational energy by NVML sampling for the two locally
served backbones, and billed token cost for the proxy-served one. Under a stated
UK grid intensity of 220\,g CO\textsubscript{2}e/kWh, a classifieds episode on
B1 corresponds to roughly 1.1--1.3\,g CO\textsubscript{2}e; the range across
modes within a cell is under 20\%.

Three things follow, and one prominently does not.

The routing result is \emph{not} an efficiency result. Since no learned policy
Pareto-beats always-cheapest, this dissertation produces no compute saving to
convert into an environmental claim, and we make none.

The oracle triage result is a measured compute reduction --- 9.5--30.6\% of cost
at zero success loss --- but it is not deployable, so it describes an
opportunity rather than a saving.

The measurement boundary excludes host CPU, cooling and facility overhead, the
browser environment's own compute, embodied emissions, and the entire remote
data centre for B0. Following \citet{gupta2022chasing} and
\citet{luccioni2022bloom}, the boundary materially determines the number, so
what we report is GPU-device operational energy and a corresponding emissions
estimate --- not a footprint. And no token reduction anywhere in this document is
converted into an energy reduction, because no hardware-independent conversion
exists \citep{fernandez2025energy}.

\section{The regime this describes}
\label{sec:regime-boundary}

Our conclusions are indexed to a regime with three named properties, and the
value of a negative result depends entirely on stating them.

\paragraph{Base success 2--36\%.} At these rates most tasks are hopeless, so
always-cheapest is close to optimal and label supply is thin. Both of the
mechanisms we identify are consequences of this number. A backbone succeeding at
60--70\% would have more labels, a smaller hopeless fraction, and a comparator
that is no longer nearly cost-optimal. \textbf{We do not claim the result
transfers there, and we would expect it not to.}

\paragraph{A cost-minimising fixed comparator.} We compare against the cheapest
fixed policy because a deployment has it for free. A study comparing against the
best fixed policy instead would report large routing wins from the same data ---
the oracle beats best-single on both axes in all eight cells. The comparator
choice is doing real work, and we consider ours the honest one; a reader who
disagrees can read the same result off Table~\ref{tab:oracle}.

\paragraph{Task-level routing.} Our decision unit is the task, not the step.
Chapter~\ref{ch:system} formulates routing per step
(Equation~\ref{eq:router}), but the labels available are task-level, because
that is where the evaluator's signal is. A step-level router might find
structure a task-level one cannot --- and would face an even worse label-supply
problem, since a step carries no outcome of its own.

\section{Limitations}
\label{sec:limitations}

These are boundaries on the conclusions, ordered by how much they would change
them.

\paragraph{Undersampling is not fully excluded as an alternative explanation.}
Our mechanism claim is that label supply is structurally constrained. A critic
can reply that 15--97 labels simply is not enough data, and that this is the
ordinary consequence of a small sample rather than a property of the regime. Two
pieces of evidence favour our reading --- three independent escape routes fail
for three different reasons (Section~\ref{sec:mechanism}), and the single
feature carrying most discriminative power is a benchmark annotation absent from
any deployment (Section~\ref{sec:matched-features}) --- but neither is decisive.
The decisive experiment is a learning curve on the most label-rich cell, plus a
demonstration that an unregularised model cannot separate the classes even in
sample. If training error is also high, the features carry no signal; if only
cross-validated error is high, the objection stands. \textbf{We did not run it,
and we state the objection as live.}

\paragraph{Replay, not live routing.} Chapter~\ref{ch:e2e}'s policies are
evaluated by crediting already-executed outcomes. This omits online overhead and
any interaction between routing and episode dynamics. The direction of the bias
favours the router, so the negative result is conservative --- but a live routed
run remains the stronger evidence, and we do not have it.

\paragraph{One benchmark family, two sites, three backbones.} Both benchmarks
are self-hosted static replicas of real sites. Live-web instability, a wider
site mix, and larger or differently trained backbones are all outside our
evidence. VisualWebArena shopping was run but excluded
(Section~\ref{sec:benchmarks}), which costs the external-validity argument a
third site it could otherwise have had.

\paragraph{B0 is not independently reproducible.} Its conditions are replicable
from archived logs but require equivalent proxy access to re-execute. No
load-bearing claim rests on B0 alone; the 0-of-8 and 8-of-8 results both hold
across backbones.

\paragraph{The noise floor is measured on one cell.} Three replicated arms exist
on \cell{cls$\cdot$B0} and one pooled comparison on \cell{wa\_red$\cdot$B1}. The
thresholds in Table~\ref{tab:floor-spread} are computed there and applied
elsewhere. When the third replicate landed it fell inside the band the first two
had established, which is reassuring but is one confirmation rather than a
distribution. We also cannot construct the six-arm repetition comparison the
ceiling would strictly require, because that needs six replicate runs of one
mode.

\paragraph{The drop-one artefact's grade stamp is an invocation artefact.} All
36 rows are \texttt{complete\_exact} and non-partial, but every row is stamped
\texttt{NON\_PAPER\_GRADE} because the producer was invoked with
\texttt{--allow-partial}. We report the stamp rather than removing it
(Appendix~\ref{app:repro}); the values support only the downgraded claim in
either case.

\section{Future work the evidence actually implies}
\label{sec:future}

\paragraph{Test the regime boundary directly.} The single most informative
follow-up is to re-run the identical protocol on a backbone with 50--70\% base
success. Our mechanism predicts that label supply relaxes, the hopeless fraction
shrinks, always-cheapest stops being near-optimal, and the winning region opens.
If routing still fails there, the mechanism claim is wrong; if it succeeds, the
regime boundary is located.

\paragraph{Settle the undersampling objection.} The learning-curve plus
in-sample-separability experiment described in Section~\ref{sec:limitations}
requires no new agent runs, only re-analysis of landed data.

\paragraph{Route on cost tiers, not modes.} Section~\ref{sec:mechanism} shows
that coarsening the target to a cost tier raises attainable agreement to
89.9--96.7\% using the same features. That is a smaller decision and a
better-posed one.

\paragraph{Report effect sizes against measured repetition floors.} The most
transferable methodological result here is negative and cheap: on
\cell{cls$\cdot$B0}, adding a different representation was worth no more than
running the same one twice. Any web-agent study reporting single-digit
percentage-point differences without a same-condition replicate is reporting a
number it cannot interpret. The replicate costs one extra run.

\paragraph{Make failure detection cheap.} The cost-first cascade attains full
oracle success at $+316$ to $+463\%$ cost, almost all of it spent confirming that
hopeless tasks are hopeless. A cheap, reliable ``this attempt has failed''
detector would convert that overspend into the prize prediction could not
claim --- and unlike the which-mode label, its supervision exists on every task.

\section{Conclusion}
\label{sec:conclusion}

Web agents pay for perception at every step of every task, and the choice of how
much to pay is normally made once and applied everywhere. This dissertation
asked whether that choice should instead be made per task, and whether it can
be.

We built a controlled representation space in which six observation modes vary
the page text, the prompt framing and the presence of an annotated image while
the task, the agent, the action space and the accounting boundary are held
fixed, and measured across eight (site, model) cells spanning two benchmarks and
three backbones.

\textbf{The opportunity is real.} A retrospective oracle gains $+3.45$ to
$+16.35\pp$ of task success over the best single mode at 1.6--35.3\% lower
cost, in every cell and on both benchmarks. \textbf{Most of its height, however, is
repetition}: one rerun arm buys 2.0--7.6\pp{} by itself, and six reruns of the
strongest single mode close to within 3.53\pp{} of the six-mode oracle --- below
the threshold we derive and apply to everyone else's effects. What survives is
the ceiling's price rather than its height: selection reaches that level for one
episode per task where repetition needs six.

\textbf{The opportunity is not reachable.} Under a fully nested protocol with a
bundle-permutation null at $B=10{,}000$, the learned triage router Pareto-beats
the trivial always-cheapest fixed policy in 0 of 8 cells --- and so does the
retrospective oracle in 7 of 8, because always-cheapest is a cost lower bound.
The negative result is therefore about the economics of the boundary, not the
capacity of the estimator.

\textbf{The reason is a supply constraint.} Routing labels are minted only by
solve events; at 2--27\% base success, 4 of 6 cells contain no trainable
which-mode classifier, and three independent routes out of the constraint fail
for three different reasons. What discriminative power the easier triage half
does show rests substantially on a feature no deployment has: the benchmark's
own difficulty annotation.

The evidence is triangulated across three backbones spanning two model families,
two benchmarks differing measurably in what their tasks specify, two deployment
classes, and --- for the load-bearing effect sizes --- a measured
same-condition repetition floor. It is scoped to self-hosted benchmark
environments at 2--36\% base success with task-level routing decisions and a
cost-minimising comparator.

What this work delivers is not a router. It is a measured account of when
representation routing can and cannot pay, and a protocol strict enough that the
answer means something: the opportunity proven to exist before it is declared
unreachable, the failure located in a named mechanism rather than left as an
implementation shortfall, and every effect size held against the noise that a
mere repetition would have produced. Those three commitments are what let a
negative result stand as evidence rather than as an absence of it --- and they
give the next system builder a way to decide, before building, whether their
regime is one in which paying for richer perception selectively could ever have
been worth it.
